% !TeX encoding = UTF-8
\SetChapterInfo{subtitle={由对偶配对、定义域与闭图像理解无界算子},keywords={对偶; 伴随; 自伴; Cayley 变换; 亏空间; 闭值域},prereq={第三章正交投影与 Riesz 表示; 第四章 Hahn--Banach; 闭图像与开映射},goals={区分转置与伴随; 把定义域纳入算子等式; 证明闭值域及可解性判据}}
\chapter{对偶算子}\label{ch:yosida-07}
\ChapterSource{按 \cite{yosida-zh,yosida-en} 第七章五节组织，2026 年 10 月 26 日补写首轮学习正文。原书位置为中译本印刷第 201--218 页、英文版第 193--208 页。实际查看英文全章扫描页及中文五节入口页；范围详见本章 COVERAGE。本稿用 $T'$ 表示线性对偶算子、$T^*$ 表示 Hilbert 伴随，内积对第一变量线性。闭值域证明的图空间步骤使用限制映射与商拓扑，回应中译本第 214 页校注指出的对偶空间识别问题。}

无界算子的公式只有连同定义域才构成算子。本章从连续泛函的延拓定义出发，经图像正交补研究闭性，再由边界条件、Cayley 变换和闭值域刻画自伴性及方程可解性。五节均包含证明、模型与附解答习题；这里的首轮写作不替代独立数学审校。

\Needspace{12\baselineskip}
\section{对偶算子}\label{sec:yosida-07-01}
\noindent\textit{原书起始页：中译本 201；英文版 193。}

本节 $X,Y$ 为 Hausdorff 局部凸空间，$X',Y'$ 表示连续线性对偶，配对写作 $\langle x,f\rangle=f(x)$，对两个变量均线性。若无特别说明，$T:D(T)\subset X\to Y$ 为线性算子，且 $D(T)$ 在 $X$ 中稠密。

\begin{definition}[对偶算子及其定义域]
令
\[
 D(T')=\{g\in Y':\text{存在 }f\in X',\quad g(Tx)=f(x)\ (x\in D(T))\},
 \qquad T'g=f.
\]
稠密性保证 $f$ 唯一：两个连续泛函在稠密子空间相等，便在全空间相等。若 $D(T)$ 不稠密，Hahn--Banach 给出非零泛函在其上为零，此时以上关系一般只能定义多值关系。
\end{definition}

\begin{proposition}[连续算子的转置]\label{prop:y7-transpose}
若 $T:X\to Y$ 连续，则 $D(T')=Y'$，且 $T'g=g\circ T$。$T'$ 对弱星拓扑 $\sigma(Y',Y)\to\sigma(X',X)$ 连续，也对强对偶拓扑 $\beta(Y',Y)\to\beta(X',X)$ 连续。若 $X,Y$ 为赋范空间，则
\[
 \|T'\|=\|T\|.
\]
\end{proposition}
\begin{proof}
复合泛函连续，故定义域为全空间。弱星连续性由 $(T'g)(x)=g(Tx)$ 直接得出。强对偶半范数为 $p_B(f)=\sup_{x\in B}|f(x)|$，其中 $B$ 遍历有界集；由 $T(B)$ 有界及 $p_B(T'g)=p_{T(B)}(g)$ 得连续性。
赋范情形有 $\|T'g\|\le\|T\|\|g\|$。反向由 Hahn--Banach 的对偶范数公式
$\|Tx\|=\sup_{\|g\|\le1}|g(Tx)|\le\|T'\|\|x\|$ 得到。
\end{proof}

\begin{proposition}[代数运算的定义域]
有界算子的线性组合满足 $(aS+bT)'=aS'+bT'$，复数标量也不取共轭。在赋范空间间，若 $T:D(T)\subset X\to Y$ 稠密定义，$S:Y\to Z$ 有界，则
\[
 (ST)'=T'S',\qquad D(T'S')=\{h\in Z':S'h\in D(T')\}.
\]
若 $S,T$ 均稠密定义、可能无界，且 $ST$ 也稠密定义，则一般只有 $T'S'\subset(ST)'$；此处复合的定义域必须逐项保留。
\end{proposition}
\begin{proof}
线性组合代入配对即可。$S$ 有界时 $D(ST)=D(T)$，故 $h(STx)=(S'h)(Tx)$ 能延拓为 $X$ 上连续泛函，当且仅当 $S'h\in D(T')$。两者无界时，对 $h\in D(T'S')$ 仍可作同样计算，但反向未必成立，因为 $D(ST)$ 可能严格小于 $D(T)$。
\end{proof}

\begin{example}[转置不保持算子强收敛]
在 $\ell^2$ 上令 $L(x_1,x_2,\ldots)=(x_2,x_3,\ldots)$，$R(x_1,x_2,\ldots)=(0,x_1,x_2,\ldots)$。有 $L^nx\to0$ 对每个 $x$ 成立。以线性对偶配对 $f_a(x)=\sum_jx_ja_j$ 识别 $\ell^2$ 与其对偶，则 $(L^n)'f_a=f_{R^na}$。非零 $a$ 满足 $\|R^na\|=\|a\|$，所以对偶算子未强收敛到零；但对每个 $x$，$((L^n)'f_a)(x)\to0$，即仍弱星收敛。
\end{example}

\begin{exercise}[矩阵与标量]\label{ex:y7-matrix}
在 $\mathbb C^n$ 上分别用双线性配对 $f_a(x)=\sum_jx_ja_j$ 和第一变量线性的标准内积表示矩阵 $A$ 的对偶算子、Hilbert 伴随。计算 $(iA)'$ 与 $(iA)^*$。
\end{exercise}
\begin{solution}
$A'$ 在 $a$ 坐标下为 $A^{\mathsf T}$；由 $(Ax,y)=(x,A^*y)$ 得 $A^*=\overline A^{\mathsf T}$。所以 $(iA)'=iA'$，$(iA)^*=-iA^*$。共轭来自 Riesz 识别的共轭线性，而不是线性对偶配对。
\end{solution}

\begin{exercise}[分布求导的转置]\label{ex:y7-distribution}
令 $D=\partial_j:\Test(\Omega)\to\Test(\Omega)$。求其对偶算子，并说明它与分布导数的关系。
\end{exercise}
\begin{solution}
测试空间上求导连续，故转置在全部 $\Dist(\Omega)$ 定义，$\langle\varphi,D'u\rangle=\langle\partial_j\varphi,u\rangle$。按分布导数定义，这等于 $-\langle\varphi,\partial_ju\rangle$，故 $D'u=-\partial_ju$。这里使用双线性分布配对，不插入复共轭。
\end{solution}

\Needspace{12\baselineskip}
\section{伴随算子}\label{sec:yosida-07-02}
\noindent\textit{原书起始页：中译本 204；英文版 195。}

本节 $H,K$ 为 Hilbert 空间，内积对第一变量线性。令 $J_H:H\to H'$ 为共轭线性等距映射 $J_Hy=(\,\cdot\,,y)$。对稠密定义 $T:D(T)\subset H\to K$，定义
\[
 D(T^*)=\{y\in K:\exists z\in H,\ (Tx,y)_K=(x,z)_H\quad\forall x\in D(T)\},
 \qquad T^*y=z.
\]
于是 $T^*=J_H^{-1}T'J_K$，其定义域也按此复合确定。等价地，$y\in D(T^*)$ 当且仅当存在 $C_y$ 使 $|(Tx,y)|\le C_y\|x\|$ 对全部 $x\in D(T)$ 成立。这里控制的是 $\|x\|$，并非图范数。

\begin{theorem}[图像、闭性和可闭性]\label{thm:y7-closure}
写 $G(T)=\{(x,Tx):x\in D(T)\}$。在 $K\oplus H$ 中有
\[
 G(T^*)=\{(-Tx,x):x\in D(T)\}^{\perp}.
\]
因此 $T^*$ 总为闭算子。$T$ 可闭当且仅当 $D(T^*)$ 在 $K$ 中稠密，此时 $\overline T=T^{**}$。特别地，稠密定义的 $T$ 闭当且仅当 $T=T^{**}$。
\end{theorem}
\begin{proof}
$(y,z)$ 与所有 $(-Tx,x)$ 正交，等价于 $-(Tx,y)+(x,z)=0$，这正是伴随定义。正交补闭，故 $T^*$ 闭。

再看图像闭包 $M=\overline{G(T)}\subset H\oplus K$。有
\[
 G(T)^\perp=\{(-T^*y,y):y\in D(T^*)\}.
\]
故 $(0,v)\in M=(G(T)^\perp)^\perp$ 当且仅当 $v\perp D(T^*)$。图像闭包为单值算子的图像，当且仅当其竖直截面只有零，即 $D(T^*)$ 稠密。这也是可闭性的充要条件。此时用同一正交补公式再取一次伴随，得到 $G(T^{**})=M$，从而 $T^{**}=\overline T$。
\end{proof}

可闭性的另一表述是：若 $x_n\in D(T)$，$x_n\to0$ 且 $Tx_n\to y$，则 $y=0$。闭图像定理只在 $D(T)=H$ 时推出有界性；稠密定义、闭、无界可以同时成立。若 $T\in\mathcal L(H,K)$，由 Riesz 表示可知 $D(T^*)=K$，且 $\|T^*\|=\|T\|$。算子包含 $S\subset T$ 总导致 $T^*\subset S^*$。

\begin{example}[一个不可闭的稠密定义算子]
在 $H=\ell^2$ 的有限支集序列空间 $c_{00}$ 上令
\[
 Tx=\Big(\sum_{n\ge1}n x_n\Big)e_1.
\]
取 $x^{(n)}=e_n/n$，则 $x^{(n)}\to0$ 而 $Tx^{(n)}=e_1$，故不可闭。由 $(Tx,y)=(\sum n x_n)\overline{y_1}$，该泛函对 $\ell^2$ 范数有界当且仅当 $y_1=0$。于是 $D(T^*)=e_1^\perp$，$T^*y=0$，其定义域果然不稠密。
\end{example}

\begin{exercise}[积分核的伴随]\label{ex:y7-kernel}
设 $k\in L^2((a,b)^2)$，$(Kf)(s)=\int_a^bk(s,t)f(t)\,\dd t$。证明 $K$ 有界并求 $K^*$。
\end{exercise}
\begin{solution}
由积分的 Cauchy--Schwarz 不等式得到 $\|Kf\|_2\le\|k\|_2\|f\|_2$。交换积分的依据为
\[
 \iint|k(s,t)f(t)\overline{g(s)}|\,\dd s\,\dd t
 \le\|k\|_2\|f\|_2\|g\|_2<\infty.
\]
因此 Fubini 定理给出
\[
 (K^*g)(t)=\int_a^b\overline{k(s,t)}g(s)\,\dd s.
\]
若 $k(s,t)=\overline{k(t,s)}$ 几乎处处，则 $K=K^*$。
\end{solution}

\begin{exercise}[图范数与闭算子]\label{ex:y7-graph}
证明闭算子 $T:D(T)\subset H\to K$ 的定义域在 $\|x\|_T=(\|x\|^2+\|Tx\|^2)^{1/2}$ 下完备。说明这为何不推出 $T$ 对原范数有界。
\end{exercise}
\begin{solution}
图范数 Cauchy 列使 $(x_n,Tx_n)$ 在 $H\oplus K$ 中 Cauchy，极限因图像闭而属于图像，故图范数完备。闭图像定理此时给出的是从图范数定义域到 $K$ 的有界性，这已由 $\|Tx\|\le\|x\|_T$ 直接成立。原范数定义域可能不完备，不能偷换两个范数。
\end{solution}

\Needspace{12\baselineskip}
\section{对称算子和自伴算子}\label{sec:yosida-07-03}
\noindent\textit{原书起始页：中译本 205；英文版 197。}

\begin{definition}[三个不同的概念]
稠密定义 $A:D(A)\subset H\to H$ 称为对称，若 $(Ax,y)=(x,Ay)$ 对 $x,y\in D(A)$ 成立，即 $A\subset A^*$；称为自伴，若 $A=A^*$，这同时要求定义域相等；称为本质自伴，若可闭且 $\overline A$ 自伴。对称算子可闭，闭包仍对称，因为 $A\subset A^*$ 给出 $\overline A=A^{**}\subset A^*$，再用 $(\overline A)^*=A^*$。
\end{definition}

\begin{example}[实乘法算子]
在 $L^2(\mathbb R)$ 上令 $M_mf=mf$，$m$ 为几乎处处有限的实可测函数，取最大定义域 $D(M_m)=\{f\in L^2:mf\in L^2\}$。这是自伴算子。
\end{example}
\begin{proof}
截断 $E_n=\{|t|\le n,\ |m(t)|\le n\}$，则 $\mathbf1_{E_n}f\to f$ 于 $L^2$，证明定义域稠密。$m$ 实保证对称。若 $g\in D(M_m^*)$ 且 $M_m^*g=h$，在 $E_n$ 上用任意 $L^2$ 探针得 $h=mg$ 几乎处处于 $E_n$。取并集得 $mg=h\in L^2$，故 $g\in D(M_m)$。于是两个定义域相等。取 $m(t)=t$ 得坐标算子，其无界性由支于 $[n,n+1]$ 的单位向量得到。
\end{proof}

\begin{example}[同一微分表达式，不同定义域]\label{exm:y7-boundary}
在 $L^2(0,1)$ 上，令 $Pf=-if'$，$D(P)=H_0^1(0,1)$。这里 $H^1$ 元素取绝对连续代表且弱导数属于 $L^2$，$H_0^1$ 等于两端迹为零的子空间。则
\[
 P^*g=-ig',\qquad D(P^*)=H^1(0,1).
\]
所以 $P$ 闭且对称，却不自伴。对 $\theta\in\mathbb R$，同一表达式在
\[
 D(P_\theta)=\{f\in H^1(0,1):f(1)=e^{i\theta}f(0)\}
\]
上定义自伴算子。
\end{example}
\begin{proof}
一维 Sobolev 事实可由积分直接看出：若 $f'=h\in L^2$（分布意义），则 $f(t)-\int_0^th(s)\,\dd s$ 的分布导数为零，故为常数。迹估计 $|f(0)|+|f(1)|\le C\|f\|_{H^1}$ 来自绝对连续公式和 Cauchy--Schwarz。零迹函数的零延拓仍在 $H^1(\mathbb R)$，经向内部缩放和光滑化可由 $C_c^\infty(0,1)$ 在 $H^1$ 中逼近，说明零迹刻画。

分部积分给出边界型
\[
 (Pf,g)-(f,-ig')=-i[f\overline g]_0^1.
\]
若 $g\in D(P^*)$，先用紧支集探针推出 $g$ 的弱导数属于 $L^2$；反向 $g\in H^1$ 时，$f$ 的零迹消去边界项。导数及迹的闭性给出 $P$ 闭。对 $P_\theta$，任意允许的端点值均能由仿射函数实现，边界型对全部 $f\in D(P_\theta)$ 消失当且仅当 $e^{i\theta}\overline{g(1)}=\overline{g(0)}$，即 $g(1)=e^{i\theta}g(0)$。于是伴随定义域恰好相等。
\end{proof}

\begin{proposition}[逆算子与满射对称算子]
若 $A$ 自伴且单射，则 $A^{-1}:R(A)\to H$ 稠密定义且自伴，逆算子未必有界。若稠密定义的对称算子 $B$ 满射，则 $B$ 自伴。
\end{proposition}
\begin{proof}
$R(A)^\perp=\ker A^*=0$ 保证逆算子稠密定义。$z\in D((A^{-1})^*)$ 且 $(A^{-1})^*z=w$，等价于 $(x,z)=(Ax,w)$ 对全部 $x\in D(A)$ 成立，亦即 $w\in D(A^*)=D(A)$ 且 $Aw=z$。故 $z\in R(A)$，$w=A^{-1}z$；反向同理，所以逆算子自伴。
对满射对称算子 $B$，有 $\ker B^*=R(B)^\perp=0$。若 $z\in D(B^*)$，选 $x\in D(B)$ 使 $Bx=B^*z$，则 $B^*(z-x)=0$，故 $z=x\in D(B)$，得到自伴。
\end{proof}

\begin{theorem}[从闭算子构造非负自伴算子]\label{thm:y7-tstar-t}
若 $T:D(T)\subset H\to K$ 闭且稠密定义，则 $T^*T$ 与 $TT^*$ 自伴且非负，且
\[
 (I+T^*T)^{-1}\in\mathcal L(H),\qquad\|(I+T^*T)^{-1}\|\le1.
\]
这里 $D(T^*T)=\{x\in D(T):Tx\in D(T^*)\}$。
\end{theorem}
\begin{proof}
闭图像的正交分解给出，对任意 $h\in H$，唯一存在 $x\in D(T),y\in D(T^*)$ 使
\[
 (h,0)=(x,Tx)+(-T^*y,y).
\]
第二分量给出 $y=-Tx$，第一分量即 $h=x+T^*Tx$。故 $I+T^*T$ 双射，且 $\|x\|\le\|h\|$，因为 $(h,x)=\|x\|^2+\|Tx\|^2$。
令 $B=(I+T^*T)^{-1}$。对 $h=(I+T^*T)x,k=(I+T^*T)z$，有 $(Bh,k)=(x,z)+(Tx,Tz)=(h,Bk)$，所以 $B$ 有界自伴且单射；$R(B)^\perp=\ker B^*=0$，说明 $D(T^*T)$ 稠密。记 $A=T^*T$。它对称且非负。若 $u\in D(A^*)$，取 $v\in D(A)$ 使 $(I+A)v=(I+A^*)u$，则 $u-v\in\ker(I+A^*)=R(I+A)^\perp=0$。故 $A=A^*$。$T^*$ 也闭且稠密定义，且 $T^{**}=T$，对 $T^*$ 重复论证即得 $TT^*$。
\end{proof}

\begin{proposition}[有界自伴算子的范数]
若 $A\in\mathcal L(H)$ 自伴，则 $\|A\|=\sup_{\|x\|=1}|(Ax,x)|$。
\end{proposition}
\begin{proof}
令右侧为 $c\le\|A\|$，$q(x)=(Ax,x)$ 为实数。由
$4t\operatorname{Re}(Ax,y)=q(x+ty)-q(x-ty)$，$t>0$，得
$|\operatorname{Re}(Ax,y)|\le c(\|x\|^2+t^2\|y\|^2)/(2t)$。取 $t=\|x\|/\|y\|$，并在复情形旋转 $y$ 的相位，使实部等于模，得到 $|(Ax,y)|\le c\|x\|\|y\|$。再对单位 $y$ 取上确界即得 $\|A\|\le c$；零向量情形显然。
\end{proof}

\begin{exercise}[亏空间检验]\label{ex:y7-deficiency}
对例 \ref{exm:y7-boundary} 的 $P$，求 $\ker(P^*-iI)$ 和 $\ker(P^*+iI)$。说明为何不能从形式表达式实对称推出自伴。
\end{exercise}
\begin{solution}
$-ig'=ig$ 给出 $g=ce^{-t}$，$-ig'=-ig$ 给出 $g=ce^t$；两者均在 $H^1(0,1)$，故两个亏空间各一维。$P$ 要求两个零端点而 $P^*$ 不要求，二者虽有相同微分表达式却不是同一算子。
\end{solution}

\begin{exercise}[实轴动量算子]\label{ex:y7-momentum}
借助第六章 Plancherel 定理，证明 $-i\,\dd/\dd t$ 在 $L^2(\mathbb R)$ 上取定义域 $H^1(\mathbb R)$ 时自伴。
\end{exercise}
\begin{solution}
令 $\mathcal U=(2\pi)^{-1/2}\FT$，则 $\mathcal U$ 酉，分布恒等式给出 $\mathcal U(-if')=\xi\mathcal Uf$。$f\in H^1$ 恰等价于 $\xi\mathcal Uf\in L^2$。故该算子为 $\mathcal U^{-1}M_\xi\mathcal U$，定义域也完整对应。实乘法算子自伴，酉变换保持伴随关系和定义域，结论成立。
\end{solution}

\Needspace{12\baselineskip}
\section{酉算子、Cayley 变换}\label{sec:yosida-07-04}
\noindent\textit{原书起始页：中译本 210；英文版 202。}

以下 $H$ 为复 Hilbert 空间。全空间上线性等距算子满足 $U^*U=I$，由极化恒等式可见保范数等价于保内积。若还满射，称为酉算子，此时 $U^{-1}=U^*$ 且 $UU^*=I$。$\ell^2$ 上右移保范数但不满射。平移 $f(t)\mapsto f(t+a)$ 及归一化 Fourier 变换均为酉算子。有界算子称为正规，若 $A^*A=AA^*$；无界情形同一等式还必须包含两个乘积的定义域相等，不能仅比较形式表达式。

\begin{theorem}[Cayley 变换及其逆]\label{thm:y7-cayley}
设 $A$ 闭、对称、稠密定义。则 $R(A\pm iI)$ 闭，且
\[
 U_A=(A-iI)(A+iI)^{-1}:R(A+iI)\longrightarrow R(A-iI)
\]
是两个闭子空间之间的满射等距映射。$I-U_A$ 单射，$R(I-U_A)=D(A)$，且
\[
 A=i(I+U_A)(I-U_A)^{-1}.
\]
反向，若 $U:M\to H$ 是定义于闭子空间 $M$ 的线性等距映射，$R(I-U)$ 稠密，则 $I-U$ 单射，以上逆公式在 $D(A)=R(I-U)$ 上定义唯一的闭对称算子，其 Cayley 变换为 $U$。
\end{theorem}
\begin{proof}
对称性使 $(Ax,x)$ 为实，故
\begin{equation}\label{eq:y7-shift-estimate}
 \|(A\pm iI)x\|^2=\|Ax\|^2+\|x\|^2.
\end{equation}
这保证单射及逆映射的界；若 $(A+iI)x_n$ 收敛，则 $x_n,Ax_n$ 均为 Cauchy，由闭性得极限仍在图像，故值域闭，负号同理。式 \eqref{eq:y7-shift-estimate} 也给出 $U_A$ 等距。写 $y=(A+iI)x$，则 $(I-U_A)y=2ix$，$(I+U_A)y=2Ax$，立即得到单射、值域与逆公式。

反向若 $Uy=y$，则对任意 $z\in M$，$(y,(I-U)z)=(y,z)-(Uy,Uz)=0$，由稠密性得 $y=0$。故可定义
$A((I-U)y)=i(I+U)y$。利用 $(Uy,Uz)=(y,z)$ 展开，得
$(i(I+U)y,(I-U)z)=((I-U)y,i(I+U)z)$，所以对称。
若 $x_n=(I-U)y_n\to x$，$Ax_n=i(I+U)y_n\to v$，则
$y_n=(x_n-iAx_n)/2$ 收敛。$M$ 闭及 $U$ 连续给出 $x=(I-U)y$、$v=i(I+U)y$，故 $A$ 闭。最后 $(A+iI)(I-U)y=2iy$、$(A-iI)(I-U)y=2iUy$，说明变换正好恢复 $U$，也证明唯一性。
\end{proof}

\begin{theorem}[自伴判据与亏空间分解]\label{thm:y7-deficiency}
对闭稠密定义的对称算子 $A$，下列等价：$A$ 自伴；$R(A+iI)=R(A-iI)=H$；$U_A$ 是 $H$ 上的酉算子。置
\[
 N_+=\ker(A^*-iI)=R(A+iI)^\perp,\qquad
 N_-=\ker(A^*+iI)=R(A-iI)^\perp.
\]
则有代数直和（不声称为 $H$ 内积正交直和）
\[
 D(A^*)=D(A)\mathbin{\dot+}N_+\mathbin{\dot+}N_-,
 \quad A^*(x_0+x_++x_-)=Ax_0+ix_+-ix_-.
\]
\end{theorem}
\begin{proof}
一般恒等式 $R(B)^\perp=\ker B^*$ 直接来自伴随定义，注意 $(A+iI)^*=A^*-iI$。若 $A=A^*$，则由式 \eqref{eq:y7-shift-estimate} 两个核均为零；两个值域既闭又稠密，故等于 $H$。反之若两者满射，任取 $y\in D(A^*)$，取 $x\in D(A)$ 使 $(A+iI)x=(A^*+iI)y$，于是 $y-x\in\ker(A^*+iI)=0$，得 $D(A^*)=D(A)$。酉判据随后由前一定理成立。

对分解，任取 $y\in D(A^*)$，在 $H=R(A+iI)\oplus N_+$ 中分解
$(A^*+iI)y=(A+iI)x_0+w_+$。令 $x_+=w_+/(2i)$，则 $x_-=y-x_0-x_+$ 属于 $N_-$，给出存在性。若 $x_0+x_++x_-=0$，施以 $A^*+iI$ 得 $(A+iI)x_0+2ix_+=0$；两项在正交子空间中，故均为零。由式 \eqref{eq:y7-shift-estimate} 得 $x_0=0$，继而 $x_+=x_-=0$。作用公式来自各子空间的定义。
\end{proof}

\begin{exercise}[乘法算子的 Cayley 变换]\label{ex:y7-cayley-mult}
求 $M_t$ 在 $L^2(\mathbb R)$ 上的 Cayley 变换，并检验逆公式的定义域。
\end{exercise}
\begin{solution}
$U$ 为乘以 $u(t)=(t-i)/(t+i)$ 的算子，$|u(t)|=1$，故为酉算子。$I-U$ 乘以 $2i/(t+i)$。其值域正是 $\{f:tf\in L^2\}$：写 $f=2ig/(t+i)$ 时 $tf\in L^2$；反之 $g=(t+i)f/(2i)\in L^2$。逆公式乘数为 $i(1+u)/(1-u)=t$，定义域也相等。
\end{solution}

\begin{exercise}[为什么两个符号都要检查]\label{ex:y7-halfline}
在 $L^2(0,\infty)$ 上取 $Af=-if'$，$D(A)=H_0^1(0,\infty)$。已知同样的弱导数论证给出 $D(A^*)=H^1(0,\infty)$。求两个亏空间，判断 $A$ 是否自伴。
\end{exercise}
\begin{solution}
$\ker(A^*-iI)$ 由 $e^{-t}$ 张成，$\ker(A^*+iI)=\{0\}$，因为 $e^t\notin L^2(0,\infty)$。$A$ 闭且对称，故 $R(A-iI)=H$，而 $R(A+iI)$ 的正交补一维。一个符号满射并不足以推出自伴；这里 $A\ne A^*$。
\end{solution}

\Needspace{12\baselineskip}
\section{闭值域定理}\label{sec:yosida-07-05}
\noindent\textit{原书起始页：中译本 214；英文版 205。}

此节 $X,Y$ 为 Banach 空间，$T:D(T)\subset X\to Y$ 闭且稠密定义。记 $N(T)=\ker T$，$R(T)=T(D(T))$。对 $M\subset X$ 写 $M^\perp=\{f\in X':f|_M=0\}$；对 $F\subset X'$ 写 ${}^\perp F=\{x\in X:f(x)=0\ \forall f\in F\}$。在 Hilbert 空间中通过 Riesz 识别后可用通常的正交补。

\begin{lemma}[有界算子的闭值域]\label{lem:y7-bounded-range}
若 $S:E\to F$ 为 Banach 空间间有界算子，则 $R(S)$ 范数闭当且仅当 $R(S')$ 范数闭。
\end{lemma}
\begin{proof}
若 $R(S)$ 闭，则 $E/N(S)\to R(S)$ 为 Banach 空间间有界双射，开映射定理给出 $\operatorname{dist}(x,N(S))\le C\|Sx\|$。对 $f\in N(S)^\perp$，令 $g_0(Sx)=f(x)$，则良定义且 $|g_0(Sx)|\le C\|f\|\|Sx\|$。Hahn--Banach 延拓给出 $g\in F'$ 且 $S'g=f$，因此 $R(S')=N(S)^\perp$ 闭。

反向设 $R(S')$ 闭。置 $Z=\overline{R(S)}$，将 $S$ 看成 $S_0:E\to Z$。限制映射 $F'\to Z'$ 满射，故 $R(S_0')=R(S')$；$S_0'$ 单射且有闭值域，逆映射定理给出 $\|g\|\le C\|S_0'g\|$（$g\in Z'$）。令 $B_E$ 为闭单位球，$K=\overline{S_0(B_E)}$。分离定理说明 $C^{-1}B_Z\subset K$：否则存在 $g$ 使某个 $\|z\|\le C^{-1}$ 满足
$\operatorname{Re}g(z)>\sup_{w\in K}\operatorname{Re}g(w)=\|S_0'g\|$，与 $|g(z)|\le C^{-1}\|g\|\le\|S_0'g\|$ 矛盾。
因此对任意余项 $r\in Z$，可选 $x$ 使 $\|x\|\le C\|r\|$ 且 $\|r-S_0x\|\le\|r\|/2$。从 $r=z$ 迭代，得到 $\sum\|x_n\|\le2C\|z\|$，余项趋零。$E$ 完备及 $S_0$ 连续保证 $S_0(\sum x_n)=z$，故 $R(S)=Z$ 闭。
\end{proof}

\begin{theorem}[Banach 闭值域定理]\label{thm:y7-closed-range}
以下四个条件等价：
\[
 \begin{aligned}
 &R(T)\text{ 在 }Y\text{ 中范数闭};\qquad
 R(T')\text{ 在 }X'\text{ 中范数闭};\\
 &R(T)={}^\perp N(T');\qquad R(T')=N(T)^\perp.
 \end{aligned}
\]
不附闭值域条件时，总有 $\overline{R(T)}={}^\perp N(T')$，但不能一般去掉闭包。
\end{theorem}
\begin{proof}
第一步将无界算子约化为有界算子，同时明确图空间的对偶。$G=G(T)$ 在 $\|(x,Tx)\|=\|x\|+\|Tx\|$ 下为 Banach 空间，令 $S:G\to Y$，$S(x,Tx)=Tx$。显然 $R(S)=R(T)$。

令 $q:X'\times Y'\to G'$ 为限制映射，$q(f,g)(x,Tx)=f(x)+g(Tx)$。Hahn--Banach 保证 $q$ 满射，它是 Banach 空间间有界满射，因而是商映射。由定义
\[
 \ker q=\{(-T'g,g):g\in D(T')\},\qquad S'g=q(0,g).
\]
于是
\[
 q^{-1}(R(S'))=\ker q+(\{0\}\times Y')=R(T')\times Y'.
\]
商拓扑说明 $R(S')$ 闭当且仅当 $R(T')$ 闭。用引理 \ref{lem:y7-bounded-range}，即得前两个条件等价。这一步把 $G'$ 处理为商空间，不能把它直接当成 $X'\times Y'$ 的子空间。

第二步，$N(T')$ 恰为 $R(T)$ 的湮灭子：若 $g$ 在 $R(T)$ 上为零，则 $g\in D(T')$ 且 $T'g=0$。Hahn--Banach 分离给出 $\overline{R(T)}={}^\perp N(T')$，所以第一、第三条件等价。

第三步，假设 $R(T)$ 闭，证明 $R(T')=N(T)^\perp$。包含 $R(T')\subset N(T)^\perp$ 直接来自配对。对反向，$N(T)$ 在 $X$ 中闭，由 $G/N(S)\to R(T)$ 的开映射定理有
\[
 \operatorname{dist}(x,N(T))+\|Tx\|\le C\|Tx\|\quad(x\in D(T)).
\]
故对 $f\in N(T)^\perp$，$g_0(Tx)=f(x)$ 良定义，且 $|g_0(Tx)|\le C\|f\|\|Tx\|$。将 $g_0$ 延拓到 $Y$ 得 $g\in Y'$，定义保证 $g\in D(T')$ 且 $T'g=f$。最后若第四条件成立，$N(T)^\perp$ 本来就是范数闭子空间，故第二条件成立，完成等价性证明。
\end{proof}

\begin{corollary}[Hilbert 空间中的可解性与估计]
若 $T:D(T)\subset H\to K$ 闭且稠密定义，则值域闭当且仅当存在 $C$ 使
\[
 \|x\|\le C\|Tx\|\quad(x\in D(T)\cap N(T)^\perp).
\]
此时 $R(T)=N(T^*)^\perp$，$R(T^*)=N(T)^\perp$。故 $Tx=f$ 可解当且仅当 $f\perp N(T^*)$，且存在唯一的 $x\perp N(T)$ 的解。
\end{corollary}
\begin{proof}
闭值域时，上述图空间开映射估计及 $\operatorname{dist}(x,N(T))=\|x\|$ 给出不等式。反向若 $Tx_n\to y$，减去 $x_n$ 的核投影不改变值域，且所得向量仍在定义域；估计使其为 Cauchy，闭性推出 $y\in R(T)$。余下等式由 Riesz 识别转译闭值域定理。
\end{proof}

当 $H=K$ 且 $\operatorname{Re}(Tx,x)\ge c\|x\|^2$（$c>0$）时，Cauchy--Schwarz 给出 $\|Tx\|\ge c\|x\|$。故 $T$ 核为零且值域闭，闭值域定理推出 $T^*$ 满射。该不等式本身不保证 $T$ 满射。

\begin{example}[稠密值域却不闭]
在 $\ell^2$ 上令 $(Tx)_n=x_n/n$，则 $T=T^*$ 为有界单射，有限支集序列均在其值域中，故值域稠密。$y=(1/n)_{n\ge1}\in\ell^2$ 不在值域，因为其唯一形式原像为 $(1,1,\ldots)\notin\ell^2$。截断 $y$ 却都在值域中。因此 $N(T^*)^\perp=\ell^2$ 严格大于 $R(T)$；也有 $\|Te_n\|=1/n$，直接否定闭值域估计。
\end{example}

\begin{exercise}[满射与下界]\label{ex:y7-surjectivity}
设 $T:D(T)\subset H\to K$ 闭且稠密定义。证明 $T$ 满射当且仅当存在 $c>0$ 使 $\|T^*y\|\ge c\|y\|$（$y\in D(T^*)$）。
\end{exercise}
\begin{solution}
满射时 $R(T)$ 闭且 $N(T^*)=R(T)^\perp=0$；闭值域定理给出 $R(T^*)$ 闭，再用前一推论的估计。反向下界使 $T^*$ 单射、值域闭（用其闭性），故 $R(T)$ 闭，又因 $N(T^*)=0$ 得 $\overline{R(T)}=K$，所以 $R(T)=K$。
\end{solution}

\begin{exercise}[周期微分方程的相容条件]\label{ex:y7-periodic}
对例 \ref{exm:y7-boundary} 的 $P_0$，证明 $R(P_0)=\{g\in L^2(0,1):\int_0^1g=0\}$，并求 $P_0f=g$ 的全部解。
\end{exercise}
\begin{solution}
积分 $-if'=g$ 得 $f(t)=c+i\int_0^tg(s)\,\dd s$。周期端点条件等价于 $\int_0^1g=0$。反向该条件下公式定义 $H^1$ 周期解；$c$ 任意，选择 $c=-i\int_0^1\int_0^tg(s)\,\dd s\,\dd t$ 得唯一零均值解。右侧是连续泛函的核，因而值域闭，也等于常数函数空间 $N(P_0^*)$ 的正交补。
\end{solution}


\Needspace{10\baselineskip}
\section*{本章参考文献与阅读衔接}
\phantomsection\addcontentsline{toc}{section}{本章参考文献与阅读衔接}
直接参考为两版吉田第七章 \cite{yosida-zh,yosida-en}。原书章末列出的 Stone、Akhiezer--Glazman、Dunford--Schwartz、Banach 以及节内 von Neumann 原始文献本轮未直接查阅。上述论证按本稿记号重新组织；Banach 闭值域定理已给出有界情形的分离与逐次逼近证明，再借助图空间商映射推广到闭算子。

第六章的形式转置决定微分表达式，本章边界条件决定 Hilbert 伴随。下一章讨论预解式和谱；式 \eqref{eq:y7-shift-estimate} 与 $T^*T$ 构造是其基础。自伴延拓分类和无界谱理论留待后续，十道习题的答案集中在书后。
